% Part: first-order-logic
% Chapter: proof-systems
% Section: natural-deduction

\documentclass[../../../include/open-logic-section]{subfiles}

\begin{document}

\iftag{FOL}
      {\olfileid[pt]{fol}{prf}{ntd}}
      {\olfileid[pt]{pl}{prf}{ntd}}

\olsection{Dedução Natural}

A dedução natural é um sistema de !!{derivation} destinado a espelhar o
raciocínio real (especialmente o tipo de raciocínio regimentado
empregado por matemáticos).  O raciocínio real procede por vários
padrões ``naturais''. Por exemplo, a prova por casos nos permite
estabelecer uma conclusão com base em uma premissa disjuntiva,
estabelecendo que a conclusão segue de qualquer um dos disjuntos. A
prova indireta nos permite estabelecer uma conclusão mostrando que sua
negação leva a uma contradição. A prova condicional estabelece uma
afirmação condicional ``se \dots então \dots'' mostrando que o
consequente segue do antecedente.  A dedução natural é uma formalização
de algumas dessas inferências naturais.  Cada um dos conectivos
lógicos e quantificadores vem com duas regras, uma de introdução e uma
de eliminação, e cada uma corresponde a um desses padrões de inferência
natural. Por exemplo, $\Intro{\lif}$ corresponde à prova condicional, e
$\Elim{\lor}$ à prova por casos.  Uma regra particularmente simples é
$\Elim{\land}$, que permite a inferência de $!A \land !B$ para~$!A$
(ou $!B$).

Uma característica que distingue a dedução natural de outros sistemas
de !!{derivation} é o uso de hipóteses. Uma !!{derivation} em dedução
natural é uma árvore de !!{formula}s.  Uma única !!{formula}
fica na raiz da árvore de !!{formula}s, e as ``folhas'' da árvore são
!!{formula}s das quais a conclusão é derivada.  Na dedução natural,
algumas !!{formula}s folha desempenham um papel dentro da
!!{derivation}, mas são ``consumidas'' quando a !!{derivation} atinge a
conclusão. Isso corresponde à prática, no raciocínio real, de introduzir
hipóteses que só permanecem em vigor por um curto período.  Por
exemplo, em uma prova por casos, assumimos a verdade de cada um dos
disjuntos; na prova condicional, assumimos a verdade do antecedente; na
prova indireta, assumimos a verdade da negação da conclusão.  Essa
maneira de introduzir hipóteses hipotéticas e depois descartá-las a
serviço de estabelecer um passo intermediário é uma marca da dedução
natural. As fórmulas nas folhas de uma !!{derivation} em dedução
natural são chamadas de hipóteses, e algumas das regras de inferência
podem descarregá-las.  Por exemplo, se temos uma !!{derivation}
de~$!B$ a partir de algumas hipóteses que incluem~$!A$, então a regra
$\Intro{\lif}$ nos permite inferir~$!A \lif !B$ e descartar qualquer
hipótese da forma~$!A$. (Para acompanhar quais hipóteses são descartadas
em quais inferências, rotulamos a inferência e as hipóteses que ela
descarta com um número.)  As hipóteses que permanecem
!!{undischarged}s no final da !!{derivation} são, juntas, suficientes
para a verdade da conclusão, de modo que uma !!{derivation} estabelece
que suas hipóteses !!{undischarged}s implicam sua conclusão.

A relação $\Gamma \Proves !A$ baseada na dedução natural vale se, e
somente se, há uma !!{derivation} em que $!A$~é a última !!{sentence}
na árvore, e toda folha que está !!{undischarged} está em~$\Gamma$.
$!A$~é um teorema na dedução natural se, e somente se, há uma
!!{derivation} em que $!A$~é a última sentença e todas as hipóteses são
!!{discharged}s. Por exemplo, aqui está uma !!{derivation} que mostra que
$\Proves (!A \land !B) \lif !A$:
\begin{prooftree}
  \AxiomC{$\Discharge{!A \land !B}{1}$}
  \RightLabel{\Elim{\land}}
  \UnaryInfC{$!A$}
  \DischargeRule{\Intro{\lif}}{1}
  \UnaryInfC{$(!A \land !B) \lif !A$}
\end{prooftree}
O rótulo~$1$ indica que a hipótese $!A \land !B$ é
!!{discharged} na inferência \Intro{\lif}.
  
Um conjunto~$\Gamma$ é inconsistente se, e somente se, $\Gamma \Proves
\lfalse$ na dedução natural.  A regra \FalseInt{} faz com que, a partir
de um conjunto inconsistente, qualquer !!{sentence} possa ser
!!{derive}da.

Sistemas de dedução natural foram desenvolvidos por Gerhard Gentzen e
Stanis\l{}aw Ja\'skowski na década de 1930, e depois por Dag Prawitz e
Frederic Fitch. Porque suas inferências espelham métodos naturais de
prova, é favorecida por filósofos. As versões desenvolvidas por Fitch
são frequentemente usadas em livros introdutórios de lógica. Na
filosofia da lógica, as regras da dedução natural às vezes foram tomadas
como dando os significados dos operadores lógicos
(``semântica da Teoria da Prova'').

\end{document}
